\section{Zero supermomentum and the relative cylinder}
\label{sec:cylinder}

\subsection{The superconformal generators}

We use the $su(1,1|1)$ generators of~\cite[App.~A]{Bulycheva:2018qcp} acting on $(\tau,\theta,\bar\theta)$, for a field of dimension $\Delta$ and R-charge $Q$:
\begin{equation}\label{eq:generators}
\begin{aligned}
L_1&=-\partial_\tau,\qquad L_0=-\tau\partial_\tau-\half\theta\partial_\theta-\half\bar\theta\partial_{\bar\theta}-\Delta,\qquad J_0=-\theta\partial_\theta+\bar\theta\partial_{\bar\theta}+Q,\\
L_{-1}&=-\tau^2\partial_\tau-\tau\theta\partial_\theta-\tau\bar\theta\partial_{\bar\theta}-2\tau\Delta-Q\,\theta\bar\theta,\\
G_{-1/2}&=\partial_{\bar\theta}-\theta\partial_\tau,\qquad G_{+1/2}=\tau\partial_{\bar\theta}-\tau\theta\partial_\tau-(2\Delta+Q)\theta-\theta\bar\theta\partial_{\bar\theta},\\
\bar G_{-1/2}&=\partial_\theta-\bar\theta\partial_\tau,\qquad \bar G_{+1/2}=\tau\partial_\theta-\tau\bar\theta\partial_\tau-(2\Delta-Q)\bar\theta+\theta\bar\theta\partial_\theta.
\end{aligned}
\end{equation}
Three remarks are needed to use these. First, the printed normalization of the R-charge is inconsistent: the coefficient $Q/2$ printed in $L_{-1}$, $G_{+1/2}$ and $\bar G_{+1/2}$ doesn't close with the constant $Q$ in $J_0$. We repair this by $Q/2\to Q$ in those three generators, as in~\eqref{eq:generators}; the alternative, $Q\to Q/2$ in $J_0$, differs only in how $Q$ is normalized. Ours is the normalization of~\cite{Fu:2016vas}, in which the supercharge has R-charge 1. Second, the printed Casimir has $-L_1L_{-1}$, with the translation on the left, while its fermionic terms have translations on the right, and it isn't central. The central operator, with translations acting first in every term, is
\begin{equation}\label{eq:casimir}
C=L_0^2-\tfrac14J_0^2-L_{-1}L_1+\half\big(G_{+1/2}\bar G_{-1/2}+\bar G_{+1/2}G_{-1/2}\big).
\end{equation}
As evidence that this ordering is what was intended, it reproduces the printed one-particle values~\cite[(A.10)--(A.12)]{Bulycheva:2018qcp} exactly when combined with the printed generators. Those values aren't constant on the multiplet, as a central element must be. With the closed algebra, $C=\Delta^2-Q^2/4$ on all components. None of this affects results at zero total charge. Third, $Q=0$ holds only for the pair. The two-point function $\ket{12}^{-2\Delta}$ is annihilated by all total generators only for $Q_1=-2\Delta$, $Q_2=+2\Delta$, in the normalization of~\eqref{eq:generators}, which is the chiral-primary condition of~\cite{Fu:2016vas}. With these corrections, a direct implementation on the Grassmann algebra reproduces the cross-ratio form of the Casimir~\cite[(4.9), (A.16)]{Bulycheva:2018qcp},
\begin{equation}\label{eq:casimir-chi}
\cC=\chi^2(1-\chi)\partial_\chi^2+\chi(1-\chi)\partial_\chi,
\end{equation}
on all odd components, with $\chi=\ket{12}\ket{34}/(\ket{14}\ket{32})$.

\subsection{Reduction}

Thong's first step is to restrict the pair Casimir to zero total momentum. The supersymmetric analogue is immediate.

\begin{proposition}\label{prop:cylinder}
Let $f$ be a function on a chiral--antichiral pair annihilated by the total $L_1$, $G_{-1/2}$ and $\bar G_{-1/2}$. Then $f=f(x)$ with $x=\ket{12}$, and
\begin{equation}\label{eq:reduced-casimir}
C\,f(x)=L_0^2f(x),\qquad L_0f(x)=-(x\partial_x+2\Delta)f(x).
\end{equation}
With $x=e^\tau n$, $n=\pm1$, and $(Wf)(\tau,n)=e^{2\Delta\tau}f(e^\tau n)$,
\begin{equation}\label{eq:WCW}
WCW^{-1}=\partial_\tau^2.
\end{equation}
\end{proposition}

\begin{proof}
Chirality leaves the coordinates $(\tau_1,\bar\theta_1;\tau_2,\theta_2)$, which is $2|2$. The three conditions remove $1|2$, leaving $1|0$, and $\ket{12}$ is the unique invariant~\cite{Bulycheva:2018qcp}. Each term of~\eqref{eq:casimir} other than $L_0^2$ and $J_0^2$ has a translation on the right, so it annihilates $f$. Also $J_0f=0$, since $\bar\theta_1\theta_2$ is neutral and $Q_1+Q_2=0$. Finally $\ket{12}$ has weight $-1$, which gives the action of $L_0$, and $W$ turns $-(x\partial_x+2\Delta)$ into $-\partial_\tau$.
\end{proof}

So \emph{the relative space has no odd directions}: the cylinder is the ordinary $\mathbb R\times S^0$, with $S^0=\{\pm1\}$. On $e^{H\tau}Y_p(n)$, with $Y_p$ of parity $\sigma_p$, the eigenvalue is $H^2$, so $\Omega=H$ and the shadow is $-H$. This derives what~\cite{Bulycheva:2018qcp} finds by computing~\eqref{eq:casimir-chi} directly.

The difference from Thong is a single term. Reordering in the same way gives $L_0^2+(1-\cN/2)L_0$ for $\cN=0,1,2$ at zero (super)momentum. With $L_0\to-h$ this is
\begin{equation}\label{eq:pattern}
h(h-d_s),\qquad d_s=1-\frac{\cN}{2},\qquad\omega=h-\frac{d_s}{2},
\end{equation}
which matches the eigenvalues $h(h-1)$~\cite{Maldacena:2016hyu}, $h(h-\half)$~\cite{Murugan:2017eto} and $h^2$~\cite{Bulycheva:2018qcp}. We also confirmed the $\cN=1$ case directly from the generators of~\cite{Murugan:2017eto}, on functions of $t_{12}-\theta_1\theta_2$. The number $d_s$ is the scaling weight of the measure on the superline with $\cN$ real odd coordinates: $d\tau\,d^{\cN}\theta$ scales with weight $-(1-\cN/2)$. So $\cN=2$ has a dimensionless measure, which is why the Casimir has no linear term. $\cN=2$ is special in another way. For $\cN=1$, the pair coordinates are $2|2$, supertranslations remove only $1|1$, and an odd relative coordinate remains. So $\cN=1$ has a true super-cylinder, and it isn't simply a smaller version of this problem.

\subsection{Components}

Expanding~\eqref{eq:bracket},
\begin{equation}\label{eq:component-relation}
f(\ket{12})=f(\tau_{12})-2\bar\theta_1\theta_2f'(\tau_{12})+\cdots,
\end{equation}
so at zero supermomentum the $b\bar b$ entry of a neutral bilocal is a derivative of the $\psi\bar\psi$ entry. In the normalization of~\eqref{eq:props} this reads $f_b=-\partial_\tau f_\psi$. The $h=0$ case is~\eqref{eq:props} itself. As a test, the matrix kernels of~\cite{Peng:2017spg} know nothing about superspace, but act on the two-component functions $(\sgn\tau\,|\tau|^{h-\delta},\ c\,|\tau|^{h-\delta-1})$ and their even partners, and on the branch $k_A$ (odd sector) or $k_S$ (even sector) their right eigenvector has exactly
\begin{equation}\label{eq:eigvec-bottom}
c=-(h-\delta)
\end{equation}
for all $h$ and $q$ tested, which is~\eqref{eq:component-relation}. The other branch has
\begin{equation}\label{eq:eigvec-top}
c=h+\delta-1,
\end{equation}
which isn't a function of $\ket{12}$ alone.

This explains the structure of the neutral problem. An ordinary centre integral imposes only $L_1=0$, and functions annihilated by $L_1$ but not by $G_{-1/2},\bar G_{-1/2}$ are still allowed. So it keeps both the bottom and the top components of each multiplet: two branches, and a matrix. Superspace keeps one function $f(x)$, one branch per parity, and one scalar kernel. We now compute that kernel.
